\subsection{FLAGS and STATUS Registers}\label{section:state-register-formats}
FLAGS and STATUS are accessed through dedicated read/write instructions.

\begin{manuallistedformatdiagram}{FLAGS Register Format}
\manualformatrow{FLAGS[15:0]}{%
\manualformatreserved{(:term:base.terms.reserved:)}{12}
\manualformatfield{Z}{1}
\manualformatfield{N}{1}
\manualformatfield{C}{1}
\manualformatfield{V}{1}
}
\end{manuallistedformatdiagram}

\manualtablecaption{FLAGS Fields}
\begin{manualtabular}{@{}>{\raggedright\arraybackslash}p{0.75in}>{\raggedright\arraybackslash}p{0.65in}>{\raggedright\arraybackslash}p{4.20in}@{}}
\toprule
\rowcolor{ManualHeaderFill}
\textbf{Field} & \textbf{Bits} & \textbf{Meaning}\\
\midrule
(:term:base.terms.reserved:) & \texttt{15..4} & reads as zero and (:term:base.terms.must_be_zero:) when written\\
\texttt{Z} & \texttt{3} & result value is zero\\
\texttt{N} & \texttt{2} & most significant result bit at the operand size\\
\texttt{C} & \texttt{1} & unsigned carry out for addition or unsigned borrow for subtraction\\
\texttt{V} & \texttt{0} & signed overflow or an explicitly reported exceptional condition\\
\bottomrule
\end{manualtabular}

(:ref:base.instructions.RDFLAGS:) and (:ref:base.instructions.WRFLAGS:) are unprivileged.
(:ref:base.instructions.RDFLAGS:) reads and zero-extends the 16-bit image; (:ref:base.instructions.WRFLAGS:) writes the low 16 bits and requires every
(:term:base.terms.reserved:) bit to be zero. Instruction-specific flag production and conditional use are defined in
\hyperref[page:flags-and-condition-codes]{Flags and Condition Codes}.

\begin{manuallistedformatdiagram}{STATUS Register Format}
\manualformatrow{STATUS[15:0]}{%
\manualformatreserved{res.}{5}
\manualformatfield{U}{1}
\manualformatfield{DEPTH}{4}
\manualformatfield{I}{1}
\manualformatfield{P}{1}
\manualformatfield{R}{1}
\manualformatfield{T}{1}
\manualformatfield{N}{1}
\manualformatfield{E}{1}
}
\end{manuallistedformatdiagram}
{\small
In the diagram, \texttt{U}, \texttt{DEPTH}, \texttt{I}, \texttt{P}, \texttt{R}, \texttt{T}, \texttt{N}, and
\texttt{E} abbreviate \texttt{UO}, \texttt{EDEPTH}, \texttt{IE}, \texttt{PM}, \texttt{RF}, \texttt{TF},
\texttt{NI}, and \texttt{EA}.
\par}

\manualtablecaption{STATUS Fields}
\begin{manualtabular}{@{}>{\raggedright\arraybackslash}p{0.75in}>{\raggedright\arraybackslash}p{0.65in}>{\raggedright\arraybackslash}p{4.20in}@{}}
\toprule
\rowcolor{ManualHeaderFill}
\textbf{Field} & \textbf{Bits} & \textbf{Meaning}\\
\midrule
(:term:base.terms.reserved:) & \texttt{15..11} & reads as zero and (:term:base.terms.must_be_zero:) when written\\
\texttt{UO} & \texttt{10} & outer active event originated in user mode\\
\texttt{EDEPTH} & \texttt{9..6} & current architectural event depth\\
\texttt{IE} & \texttt{5} & permits maskable-interrupt admission when the other admission conditions hold\\
\texttt{PM} & \texttt{4} & selects user mode when zero and supervisor mode when one\\
\texttt{RF} & \texttt{3} & one-shot suppression of (:ref:base.events.EXCEPTION.DEBUG_TRACE:)\\
\texttt{TF} & \texttt{2} & requests post-success (:ref:base.events.EXCEPTION.DEBUG_TRACE:) for a (:term:base.terms.trace_unit:) when \texttt{RF} is clear\\
\texttt{NI} & \texttt{1} & inhibits NMI admission while preserving the NMI pending\\
\texttt{EA} & \texttt{0} & hardware-managed architectural event-active state\\
\bottomrule
\end{manualtabular}

In user mode, \texttt{PM}, \texttt{EA}, \texttt{EDEPTH}, and \texttt{UO} are always observed as zero.

(:ref:base.instructions.RDSTATUS:) is unprivileged; (:ref:base.instructions.WRSTATUS:) is
supervisor-only. (:ref:base.instructions.WRSTATUS:) atomically updates \texttt{IE}, \texttt{NI}, \texttt{TF}, and \texttt{RF};
(:term:base.terms.reserved:) bits (:term:base.terms.must_be_zero:), and supplied \texttt{PM}, \texttt{EA}, \texttt{EDEPTH}, and \texttt{UO} must equal
their current values. Privilege
transitions are defined in the \hyperref[page:privileged-execution-model]{Privileged Execution Model}; event
admission, tracing, and event-active transitions are defined in the
\hyperref[page:architectural-event-processing-model]{Architectural Event Processing Model}.
